\section{Case Study: Dependency-Free FlashAttention for Hopper and Blackwell}
\label{sec:kernels}
% Design details, the limit analysis, the full figure, and per-scenario numbers are in
% Appendix~\ref{app:kernels} and Appendix~\ref{app:per-scenario}.

\paragraph{Task.} We asked \sys to produce a FlashAttention forward kernel for Hopper and one for
Blackwell, each a single CUDA file with no CUTLASS, CuTe, cuBLAS, or cuDNN, exposing the operator
schema vLLM uses for FlashAttention-3 so that the result drops in for the upstream kernel and is
timed against it from identical inputs. The scenario suites mirror the upstream kernels' own test
grids and add production shapes from SOL-ExecBench \cite{lin2026solexecbench}: 182 scenarios per
architecture spanning bf16, fp16, and fp8; paged and contiguous KV caches; MHA, GQA, and MQA; causal,
sliding-window, softcapped, and attention-sink masking; ragged batches; prefill, chunked prefill,
and split-KV decode; and unequal query and value head dimensions. The breadth is deliberate, since kernels scored on a few
shapes report speedups that shrink on the production path \cite{oliaro2026fastkernels}. The goal was correctness on every scenario and at least parity with upstream in geometric mean. The seed
in both cases was a slow but general kernel (a scalar one-warp-per-query kernel on Hopper, a
one-warpgroup dense-only skeleton on Blackwell), and anything the fast paths did not yet cover fell
back to it, so correctness held throughout. The upstream kernels, FA3 as shipped in vLLM and the
CuteDSL FA4 \cite{shah2024flashattention3,zadouri2026flashattention4}, were the correctness oracle
and the yardstick the goal was stated against. The dependency-free constraint means the layouts,
descriptors, pipelines, and register allocation that CUTLASS and CuteDSL derive for the upstream
kernels had to be written by hand \cite{nvidia2023cutlass}.

\paragraph{What the agent built.} On Hopper, \fathree (3.6K lines) is a warp-specialized kernel in
raw CUDA and PTX: producer and consumer warpgroups connected by a TMA and mbarrier ring,
asynchronous \texttt{wgmma} MMAs with \texttt{setmaxnreg} register rebalancing, persistent
scheduling with inter-warpgroup ping-pong, split-KV decode with a combine pass, and a native fp8
path. Its one remaining limit sits where a library's allocator would otherwise be: at head dimension
256 the hand-written swizzle composes only at a few tile widths and the output accumulator caps
occupancy at one CTA per SM, so that scenario runs at 0.94$\times$. On Blackwell, \fafour (10.2K
lines) is a family with host-side routing, a structure the questions pushed the agent toward after a
single kernel had stalled (\S\ref{sec:eval-cases}): a persistent warp-specialized collective for the
dense envelope (two stacked 128-row query halves, two softmax warpgroups in an ordered ping-pong, a
correction warpgroup rescaling the output in tensor memory, one MMA warp driving \texttt{tcgen05}),
variants for unequal head dimensions and paged prefill, and a two-tile kernel for small, ragged, and
decode shapes. Several decisive gains were outside the kernel bodies: a synchronous host copy
costing about 38~$\mu$s per call, a stale routing gate, a TMA gather issued from divergent lanes, and
ring state the compiler had placed in local memory. fp8 inputs and dense head dimensions of 192 and
above are not yet built on Blackwell (28 of 182 scenarios).

\begin{table}[t]
\centering\small
\caption{Case-study results. $n$ counts the scenarios the upstream kernel also runs (paired).
``Upstream'' is FA3 on Hopper and FA4 on Blackwell; ``open-source'' is, per scenario, the faster of
upstream and FlashInfer; ``available'' also includes cuDNN's fused attention. Speedups are geometric
means of wall-clock ratios from byte-identical inputs timed back to back on one GPU
(Appendix~\ref{app:bench}).}
\label{tab:casestudy}
\setlength{\tabcolsep}{3pt}\footnotesize
\begin{tabular}{lrrrrr}
\toprule
Set & $n$ & Correct & \multicolumn{3}{c}{Speedup vs.\ fastest \ldots} \\
\cmidrule(lr){4-6}
 & & & upstream & open-source & available \\
\midrule
\multicolumn{6}{l}{\emph{Hopper (H200), upstream = FA3}} \\
attntest   & 64  & 64/64   & 1.22$\times$ & 1.22$\times$ & 1.14$\times$ \\
edge       & 31  & 31/31   & 1.55$\times$ & 1.55$\times$ & 1.31$\times$ \\
claim      & 77  & 77/77   & 1.01$\times$ & 1.01$\times$ & 0.99$\times$ \\
real       & 10  & 10/10   & 1.10$\times$ & 1.10$\times$ & 1.08$\times$ \\
all        & 182 & 182/182 & 1.16$\times$ & 1.16$\times$ & 1.10$\times$ \\
\midrule
\multicolumn{6}{l}{\emph{Blackwell (B200), upstream = FA4}} \\
attntest   & 33  & 52/52   & 1.44$\times$ & 1.42$\times$ & 1.09$\times$ \\
edge       & 28  & 31/31   & 2.02$\times$ & 1.93$\times$ & 1.42$\times$ \\
claim      & 55  & 61/61   & 1.30$\times$ & 1.20$\times$ & 0.92$\times$ \\
real       & 8   & 10/10   & 1.10$\times$ & 1.10$\times$ & 1.02$\times$ \\
all        & 124 & 154/154 & 1.46$\times$ & 1.39$\times$ & 1.07$\times$ \\
\bottomrule
\end{tabular}
\vspace{2pt}

\raggedright\footnotesize ``Correct'' is over runnable scenarios. At or above parity vs.\ upstream /
open-source / available: Hopper 135 / 135 / 121 of 182; Blackwell 98 / 94 / 63 of 124. On Hopper,
cuDNN supports 145 scenarios and \fathree is faster on 111; FlashInfer supports 32 and \fathree is
faster on all 32.
\end{table}

\paragraph{Results.} Table~\ref{tab:casestudy} summarizes the outcome; Appendix~\ref{app:kernels}
plots every scenario. Both kernels are correct on every scenario they run, and both beat upstream in
geometric mean. \fathree runs at 1.16$\times$ FA3 with 135 of 182 scenarios at or above parity; the
wins are concentrated in ragged prefill, paged and speculative decode, and the masking grid, while
the paper-style sweeps of the \texttt{claim} set are at parity and dense fp8 prefill is the one clear
hole at 0.67$\times$. \fafour runs at 1.46$\times$ FA4 across the 124 scenarios FA4 also runs, with
98 at or above parity and a worst case of 0.78$\times$. The headline shrinks as the baseline
strengthens: cuDNN is the fastest baseline on 54 of 182 Hopper scenarios and on 92 of 124 Blackwell
ones, where it wins dense prefill and batched decode outright, so the kernels are 1.10$\times$ and
1.07$\times$ the fastest available implementation, and at or above it on production shapes. Three caveats are documented in Appendix~\ref{app:kernels}: the \texttt{edge}
sets consist of short kernels, so part of their 1.5--2$\times$ is upstream dispatch overhead; the
largest \texttt{claim} wins are fp8-KV decode shapes on which FA4 is slow, and excluding them moves
the Blackwell all-set geomeans to 1.37, 1.32, and 1.08$\times$; and one paged-KV fast path depends on
physically consecutive pages. Every Blackwell scenario below parity has a measured cause: dense long
context spends identical tensor work over a third more cycles because of twice as many barrier waits
per block, batched decode pays a cold start per CTA, and head dimension 96 runs on a padded 128-wide
path. On Hopper the losses are fp8 dense prefill, softcapped prefill, head dimension 96, and long
sliding windows (0.34--0.79$\times$); these have not been profiled to the same depth.

\paragraph{Release.} We release both kernels under Apache~2.0 with the evaluator, the PTX probes used
to validate each Blackwell primitive in isolation, the search scripts (\S\ref{sec:workflow-evo}),
the question bank and skill, and the complete agent traces of both approaches. The kernels are meant as seeds, with their residual
gaps measured, so that the next agent starts from a known frontier.
